% Generated draft: requires mathematical and source review.
\documentclass[11pt]{article}
\usepackage[margin=25mm]{geometry}
\usepackage{fontspec}
\setmainfont{Latin Modern Roman}
\usepackage{amsmath,amssymb,mathtools,unicode-math}
\setmathfont{Latin Modern Math}
\usepackage{enumitem,needspace,xurl,hyperref}
\setlength{\parindent}{0pt}
\setlength{\parskip}{5pt}
\setlength{\emergencystretch}{3em}
\newcommand{\N}{\mathbb{N}}
\newcommand{\Z}{\mathbb{Z}}
\newcommand{\Q}{\mathbb{Q}}
\newcommand{\R}{\mathbb{R}}
\newcommand{\C}{\mathbb{C}}
\newcommand{\F}{\mathbb{F}}
\newcommand{\A}{\mathbb{A}}
\newcommand{\PP}{\mathbb{P}}
\newcommand{\E}{\mathbb{E}}
\newcommand{\eps}{\varepsilon}
\newcommand{\dd}{\,\mathrm{d}}
\begin{document}
\section*{UnsolvedMath ID: 800208}
\textbf{Source problem code:} catalogue entry
\subsection{Definitions and mathematical statement}
Let \(M\) and \(N\) be strongly pseudoconvex CR manifolds of real dimension \(2n+1\), with contact distributions \(H_M\subset TM\) and \(H_N\subset TN\), and CR complex structures \(J_M\) and \(J_N\) on these distributions. Here \(n\geq 1\) is the usual positive CR dimension. Let \(f:M\to N\) be a fixed diffeomorphism satisfying
\[
\dd f_x(H_{M,x})=H_{N,f(x)}\qquad(x\in M).
\]
The source calls such a map quasiconformal when it distorts the CR structures boundedly, and assumes that \(f\) is quasiconformal.

A Hermitian metric on a contact bundle is understood here through its associated positive definite real inner product \(g\), satisfying \(g(Ju,Jv)=g(u,v)\). Given Hermitian metrics \(g\) on \(H_M\) and \(h\) on \(H_N\), put \(A_x=\dd f_x|_{H_{M,x}}\). The ratio of the longest to the shortest semiaxis of the image of the \(g_x\)-unit sphere is
\[
K_{g,h}(x)=
\frac{\displaystyle\max_{v\in H_{M,x},\ g_x(v,v)=1}
\sqrt{h_{f(x)}(A_xv,A_xv)}}
{\displaystyle\min_{v\in H_{M,x},\ g_x(v,v)=1}
\sqrt{h_{f(x)}(A_xv,A_xv)}}.
\]
This gives a precise interpretation of the source's sphere-to-ellipsoid measurement. The source explicitly states that the measurement depends on the chosen Hermitian metrics when \(n>1\).

\textbf{Source question.} For the fixed quasiconformal diffeomorphism \(f\), can Hermitian metrics on the two contact bundles be chosen so that its distortion is minimal?

\textbf{Unresolved optimization convention.} The source does not specify whether minimality is pointwise or refers to a global distortion bound, nor does it specify the admissible metric class or its regularity. For example, if both metrics may vary over all smooth Hermitian metrics and distortion means the global supremum
\[
\mathcal K_f(g,h)=\sup_{x\in M}K_{g,h}(x)\in[1,+\infty],
\]
then a possible precise formulation is to ask whether there exist smooth Hermitian metrics \(g_*,h_*\) such that
\[
\mathcal K_f(g_*,h_*)=
\inf_{\substack{g\text{ smooth Hermitian on }H_M\\
                 h\text{ smooth Hermitian on }H_N}}
\mathcal K_f(g,h).
\]
This displayed optimization is conditional on those conventions; the supplied record does not establish that they are the intended ones. No compactness, boundary, metric normalization, or uniform comparison requirements are supplied.
\subsection{Short English statement}
For a fixed quasiconformal diffeomorphism between strongly pseudoconvex CR manifolds, can one choose Hermitian metrics on their contact bundles that minimize its distortion? The source notes metric dependence when \(n>1\), but leaves the meaning of minimum and the allowable metric choices unspecified.
\subsection{Sources}
\begin{itemize}
\item Local UnsolvedMath problem registry, permanent ID \(800208\), exact source problem code \texttt{catalogue entry}; title: \emph{Several Complex Variables Problem List --- Problem 4}. The supplied registry labels the problem open; that label is not an independently verified assessment.
\item Buzzard, Fornaess, and Krantz, \emph{The Several Complex Variables Problem List} (1995), Problem 4. Supplied source URL: \url{https://arxiv.org/abs/math/9509202}. The supplied background records access on 2026-10-16 and extraction by \texttt{arxiv\_latex}; these are supplied provenance statements, not a claim of independent source access.
\item The extracted problem is attributed to Puqi Tang, \texttt{tang@math.purdue.edu}, and carries the qualifications ``(New)'' and ``(See Section 11.3)''. The registry's \texttt{proposed\_by} field names Buzzard, Fornaess, and Krantz and gives the proposed year as 1995. The referenced Section 11.3 is not included in the supplied record.
\end{itemize}
\end{document}
